% META: {"id": "TCOMPL-TEMPS-ATTENTE-QCM", "chapitre": "TCOMPL-TEMPS-ATTENTE", "type_objet": "qcm", "genere_depuis": "chapitres/TCOMPL-TEMPS-ATTENTE/qcm/TCOMPL-TEMPS-ATTENTE-QCM.json", "status": "generated"}
% Fichier genere par scripts/build_qcm_tex.py — ne pas editer a la main.

\section*{\textcolor{chapcolor}{\MakeUppercase{Loi géométrique et temps d'attente}}}

\begin{center}
\textit{Pour chaque question, une seule réponse est exacte.}
\end{center}

\begin{enumerate}

\item[] \textbf{Capacité C1}

\item \textbf{[Q1]} Une variable aleatoire X suivant une loi géométrique de paramètre p a pour espérance :
  \begin{enumerate}[label=\Alph*.]
    \item $\frac{1}{p}$
    \item $\frac{1-p}{p^2}$
    \item $p$
    \item $1-p$
  \end{enumerate}

\item[] \textbf{Capacité C2}

\item \textbf{[Q2]} Pour $k \ge 1$, la probabilité $P(X = k)$ d'une loi géométrique $G(p)$ vaut :
  \begin{enumerate}[label=\Alph*.]
    \item $(1-p)^{k-1} p$
    \item $p (1-p)^k$
    \item $p^k (1-p)$
    \item $\binom{n}{k} p^k (1-p)^{n-k}$
  \end{enumerate}

\item[] \textbf{Capacité C3}

\item \textbf{[Q3]} La propriété de perte de memoire de la loi géométrique s'ecrit $P_{X > n}(X > n+k) =$
  \begin{enumerate}[label=\Alph*.]
    \item $P(X > n)$
    \item $P(X > k)$
    \item $P(X = k)$
    \item $1 - P(X > k)$
  \end{enumerate}

\item[] \textbf{Capacité C2}

\item \textbf{[Q4]} Pour une loi géométrique $G(p)$, $P(X > k) =$
  \begin{enumerate}[label=\Alph*.]
    \item $p^k$
    \item $1 - p^k$
    \item $(1-p)^k$
    \item 1 - (1-p)\textasciicircum{}k
  \end{enumerate}

\item[] \textbf{Capacité C1}

\item \textbf{[Q5]} Si $p = 0{,}2$, le temps d'attente moyen du premier succès est :
  \begin{enumerate}[label=\Alph*.]
    \item 2 d'essais
    \item 4 d'essais
    \item 10 d'essais
    \item 5 d'essais
  \end{enumerate}

\item[] \textbf{Capacité C4}

\item \textbf{[Q6]} La loi exponentielle de paramètre $\lambda>0$ est dite « sans mémoire » parce que, pour tous $s,t\geq0$ :
  \begin{enumerate}[label=\Alph*.]
    \item $P(X>s+t)=P(X>s)+P(X>t)$
    \item $P_{X>s}(X>s+t)=P(X>t)$
    \item $P(X>t)=1-e^{-\lambda t}$
    \item $E(X)=\lambda$
  \end{enumerate}

\item[] \textbf{Capacité C5}

\item \textbf{[Q7]} Pour que la fonction $f$ définie par $f(x)=k\,e^{-2x}$ sur $[0;+\infty[$ et nulle ailleurs soit une densité de probabilité, il faut :
  \begin{enumerate}[label=\Alph*.]
    \item $k=1$
    \item $k=\tfrac12$
    \item $k=2$
    \item n'importe quel $k>0$
  \end{enumerate}

\item[] \textbf{Capacité C6}

\item \textbf{[Q8]} Pour $X$ de loi exponentielle de paramètre $\lambda$, l'espérance s'écrit :
  \begin{enumerate}[label=\Alph*.]
    \item $\displaystyle\int_0^{+\infty} x\,e^{-\lambda x}\,dx$
    \item $\displaystyle\int_{-\infty}^{+\infty} \lambda x\,e^{-\lambda x}\,dx$
    \item $\displaystyle\int_0^{+\infty} \lambda e^{-\lambda x}\,dx$
    \item $\displaystyle\int_0^{+\infty} x\,\lambda e^{-\lambda x}\,dx$, qui vaut $\tfrac1\lambda$
  \end{enumerate}

\item[] \textbf{Capacité C7}

\item \textbf{[Q9]} Si $X$ suit la loi uniforme sur $[a;b]$ (avec $a<b$), alors :
  \begin{enumerate}[label=\Alph*.]
    \item $E(X)=\tfrac{a+b}2$ et $V(X)=\tfrac{(b-a)^2}{12}$
    \item $E(X)=\tfrac{b-a}2$ et $V(X)=\tfrac{(b-a)^2}{12}$
    \item $E(X)=\tfrac{a+b}2$ et $V(X)=\tfrac{b-a}{12}$
    \item sa densité vaut $\tfrac1b$ sur $[a;b]$
  \end{enumerate}

\end{enumerate}

% NEXUS-QCM-TEACHER-ONLY-BEGIN
\ifnxVersionProfesseur
\clearpage
\section*{Clé de correction — réservée au professeur}

\begin{center}
\begin{tabular}{lll}
\hline
Question & Capacité & Réponse exacte \\
\hline
Q1 & C1 & \textbf{A} \\
Q2 & C2 & \textbf{A} \\
Q3 & C3 & \textbf{B} \\
Q4 & C2 & \textbf{C} \\
Q5 & C1 & \textbf{D} \\
Q6 & C4 & \textbf{B} \\
Q7 & C5 & \textbf{C} \\
Q8 & C6 & \textbf{D} \\
Q9 & C7 & \textbf{A} \\
\hline
\end{tabular}
\end{center}

\subsection*{Diagnostic des réponses erronées}

\begin{itemize}
  \item \textbf{Q1}
  \begin{itemize}
    \item \textbf{B} — Formule de la variance. Avec la convention X=nombre d'essais jusqu'au premier succès, l'espérance de G(p) vaut 1/p. \emph{Renvoi : C1.}
    \item \textbf{C} — Paramètre p. Avec la convention X=nombre d'essais jusqu'au premier succès, l'espérance de G(p) vaut 1/p. \emph{Renvoi : C1.}
    \item \textbf{D} — Probabilité d'echec. Avec la convention X=nombre d'essais jusqu'au premier succès, l'espérance de G(p) vaut 1/p. \emph{Renvoi : C1.}
  \end{itemize}
  \item \textbf{Q2}
  \begin{itemize}
    \item \textbf{B} — Erreur d'exposant k au lieu de k-1. Avoir le premier succès au rang k exige k-1 échecs puis un succès, donc P(X=k)=(1-p)\textasciicircum{}(k-1)p. \emph{Renvoi : C2.}
    \item \textbf{C} — Inversion des roles de succès et echec. Avoir le premier succès au rang k exige k-1 échecs puis un succès, donc P(X=k)=(1-p)\textasciicircum{}(k-1)p. \emph{Renvoi : C2.}
    \item \textbf{D} — Formule de la loi binomiale. Avoir le premier succès au rang k exige k-1 échecs puis un succès, donc P(X=k)=(1-p)\textasciicircum{}(k-1)p. \emph{Renvoi : C2.}
  \end{itemize}
  \item \textbf{Q3}
  \begin{itemize}
    \item \textbf{A} — Temps n au lieu de k. L'absence de mémoire donne P(X>n+k | X>n)=P(X>k). \emph{Renvoi : C3.}
    \item \textbf{C} — Égalité au lieu d'une superiocrite stricte. L'absence de mémoire donne P(X>n+k | X>n)=P(X>k). \emph{Renvoi : C3.}
    \item \textbf{D} — Événement contraire. L'absence de mémoire donne P(X>n+k | X>n)=P(X>k). \emph{Renvoi : C3.}
  \end{itemize}
  \item \textbf{Q4}
  \begin{itemize}
    \item \textbf{A} — Utilisation de p au lieu de (1-p). L'événement X>k signifie k échecs successifs, donc P(X>k)=(1-p)\textasciicircum{}k. \emph{Renvoi : C2.}
    \item \textbf{B} — Formule fausse. L'événement X>k signifie k échecs successifs, donc P(X>k)=(1-p)\textasciicircum{}k. \emph{Renvoi : C2.}
    \item \textbf{D} — Probabilité P(X <= k). L'événement X>k signifie k échecs successifs, donc P(X>k)=(1-p)\textasciicircum{}k. \emph{Renvoi : C2.}
  \end{itemize}
  \item \textbf{Q5}
  \begin{itemize}
    \item \textbf{A} — Calcul 1/0.5. Pour p=0,2, le temps moyen vaut E(X)=1/p=5 essais. \emph{Renvoi : C1.}
    \item \textbf{B} — Calcul (1-p)/p = 0.8/0.2 = 4. Pour p=0,2, le temps moyen vaut E(X)=1/p=5 essais. \emph{Renvoi : C1.}
    \item \textbf{C} — Calcul 1/(0.2\textasciicircum{}2). Pour p=0,2, le temps moyen vaut E(X)=1/p=5 essais. \emph{Renvoi : C1.}
  \end{itemize}
  \item \textbf{Q6}
  \begin{itemize}
    \item \textbf{A} — Les probabilités de survie se multiplient : $e^{-\lambda(s+t)}=e^{-\lambda s}e^{-\lambda t}$. C'est ce produit qui, une fois divisé par $P(X>s)$, donne l'absence de mémoire. \emph{Renvoi : C4.}
    \item \textbf{C} — $1-e^{-\lambda t}$ est la fonction de répartition $P(X\leq t)$ ; la probabilité de survie vaut $P(X>t)=e^{-\lambda t}$. Et une formule de répartition n'exprime pas, à elle seule, l'absence de mémoire. \emph{Renvoi : C4.}
    \item \textbf{D} — L'espérance vaut $\tfrac1\lambda$ et non $\lambda$ ; surtout, une espérance ne dit rien du conditionnement par $\{X>s\}$. \emph{Renvoi : C4.}
  \end{itemize}
  \item \textbf{Q7}
  \begin{itemize}
    \item \textbf{A} — Avec $k=1$, $\int_0^{+\infty}e^{-2x}\,dx=\tfrac12\neq1$. La condition $\int f=1$ impose $\tfrac k2=1$, donc $k=2$. \emph{Renvoi : C5.}
    \item \textbf{B} — Avec $k=\tfrac12$ l'intégrale vaut $\tfrac14$. Le facteur cherché est l'inverse de $\int_0^{+\infty}e^{-2x}\,dx=\tfrac12$, soit $k=2$. \emph{Renvoi : C5.}
    \item \textbf{D} — La positivité ne suffit pas : une densité doit aussi avoir une intégrale égale à $1$ sur $\mathbb{R}$, ce qui fixe $k=2$. \emph{Renvoi : C5.}
  \end{itemize}
  \item \textbf{Q8}
  \begin{itemize}
    \item \textbf{A} — Le facteur $\lambda$ de la densité a été oublié : cette intégrale vaut $\tfrac1{\lambda^2}$, pas $\tfrac1\lambda$. L'espérance intègre $x$ contre la densité $\lambda e^{-\lambda x}$. \emph{Renvoi : C6.}
    \item \textbf{B} — La densité est nulle sur $]-\infty;0[$ : l'intégrale doit partir de $0$. Écrite de $-\infty$ à $+\infty$ avec $\lambda x e^{-\lambda x}$, elle diverge. \emph{Renvoi : C6.}
    \item \textbf{C} — C'est l'intégrale de la densité seule : elle vaut $1$, la masse totale de probabilité. L'espérance pondère chaque valeur $x$ par cette densité. \emph{Renvoi : C6.}
  \end{itemize}
  \item \textbf{Q9}
  \begin{itemize}
    \item \textbf{B} — $\tfrac{b-a}2$ est la demi-longueur de l'intervalle, pas son milieu : sur $[2;6]$ elle vaudrait $2$, alors que le centre est $4=\tfrac{2+6}2$. \emph{Renvoi : C7.}
    \item \textbf{C} — Une variance est homogène au carré de la variable : elle vaut $\tfrac{(b-a)^2}{12}$. La forme $\tfrac{b-a}{12}$ changerait de valeur avec l'unité choisie. \emph{Renvoi : C7.}
    \item \textbf{D} — La densité est constante égale à $\tfrac1{b-a}$ sur $[a;b]$, seule valeur donnant une intégrale égale à $1$ ; $\tfrac1b$ ne convient que si $a=0$. \emph{Renvoi : C7.}
  \end{itemize}
\end{itemize}
\fi
% NEXUS-QCM-TEACHER-ONLY-END
